\documentclass[../main.tex]{subfiles}
\begin{document}
\setchapterstyle{kao}
\setchapterpreamble[u]{\margintoc}
\setchapterimage[6.5cm]{Images/intro.jpg}
\chapter[Introduction]{Introduction\footnotemark[0]}
\labch{intro}
\section{Dimensional Analysis and Energy Scales}
We start our discussion about the Standard Model by introducing some fundamental concepts about dimensional analysis. For our purposes, we work with natural units, i.e. $\hbar=c=1$. It follows that:
\begin{itemize}
    \item $[m]=[E]$ since $E=mc^2$ 
    \item $[t]=[E^{-1}]$ since $\Delta E\Delta t\gtrsim\hbar$
    \item $P^\mu=(E,\Vec{p})$ and we have that $[x]=[E^{-1}]$
    \item the Lagrangian density $\pazocal{L}$ has dimension $[\pazocal{L}]=[E^4]$ because it has to be $[S]=\left[\int d^4x\pazocal{L}\right]=0$
    \item The fields we work with have dimensions $[\phi]=[E]$, $[\partial_\mu\phi]=[E]$, $[\Psi]=[E^{3/2}]$, $[A_\mu]=[E]$
\end{itemize}
For the couplings, we observe that $e\bar{\Psi}\gamma^\mu A_\mu\Psi$ must have dimension $E^4$ therefore it has to be $[e]=[E^0]$. Same reasoning applies for $G_F\bar{\Psi}\Psi\bar{\Psi}\Psi$ and we obtain $[G_F]=[E^{-2}]$.

The typical energy scale we are going to use is GeV, while the masses we are interested in are:
\[
\left\{
\begin{aligned}
&\text{Leptons: } &&m_e=0.511\text{ MeV} \quad &&m_\mu=105\text{ MeV} \quad &&m_\tau=1.77\text{ GeV}\\
&\text{Light quarks: } &&m_u=2.5\text{ MeV} \quad &&m_d=5.0\text{ MeV} \quad &&m_s=0.105\text{ GeV}\\
&\text{Heavy quarks: } &&m_c=1.27\text{GeV} \quad &&m_b=4.20\text{GeV} \quad &&m_t=173\text{GeV}
\end{aligned}
\right.
%Grazie a Giorgio che mi ha fatto allineare tutto per bene (giovedì 20 ottobre 2022, nell'Ufficio)
\]
In quantum chromodynamics (QCD) the QCD scale is given by the quantity $\Lambda_{QCD}$. Its value is $\Lambda_{QCD}=332\pm17$ MeV when the energy involved in the process allows only to produce the light quarks, i.e. energies below 1.27 GeV. At higher energy, $\Lambda_{QCD}=210\pm14$ MeV, above the bottom quark mass. This gives mass to elementary fields:
\[
\left\{
\begin{aligned}
&m_p=938\text{ MeV}=g_p\Lambda_{QCD} \quad m_n=939\text{ MeV}=g_n\Lambda_{QCD}\\
&m_{\pi^\pm}=139\text{ MeV}, m_{\pi^0}=135\text{ MeV}\propto\sqrt{(m_u+m_d)\Lambda_{QCD}}
\end{aligned}
\right.
\Rightarrow \text{we have a different scaling}
\]
Lastly, we give the mass of the two gauge bosons $W$ and $Z$ and of the Higgs boson $h$:\marginnote{Remember that the other two gauge bosons $\gamma$ and $g$ are massless.}
\[
m_W\simeq80\text{ GeV} \quad m_Z\simeq90\text{ GeV} \quad m_h\simeq125\text{ GeV}
\]
\subsection{Muon Decay}
With the energy scales briefly discussed before, we want to compute the muon decay width and its lifetime.\\
For a decay of the type $A\xrightarrow[]{}B_1+B_2+B_3+\dots$ we know that
\[
d\Gamma=\frac{1}{2m_A}\left(\prod_f\frac{d^3p_A}{(2\pi)^3}\frac{1}{2E_f}\right)(2\pi)^4\delta^4(p_A-\{p_f\})|\pazocal{M}|^2
\]
We are now interested in $\mu^-\xrightarrow[]{}e^-\bar{\nu}_e\nu_\mu$
\marginnote[-15mm]{\includegraphics[]{Images/muondecay.pdf}} and what we want to compute is
\[
\pazocal{M}\{\mu^-(P_\mu,a)\xrightarrow[]{}e^-(P_e,b)+\bar{\nu}_e(\underset{\mathclap{\tikz \node {$\uparrow$} node [below=1ex] {\footnotesize $P_1$};}}{P_{\bar{\nu}_e}},c)+\nu_\mu(\underset{\mathclap{\tikz \node {$\uparrow$} node [below=1ex] {\footnotesize $P_2$};}}{P_{\nu_\mu}},d)\}
\]
The gauge bosons are massive, so we need to fix the gauge: in the unitary gauge, we have that the propagator is given by 
\[
-i\frac{g^{\alpha\beta}-k^\alpha k^\beta/m_W^2}{K^2-m^2_W+i\varepsilon}
\]
with $K^\alpha=(P_\mu)^\alpha-(P_2)^\alpha$. We need an additional Feynman rule for this diagram: for each vertex in which there is an interaction with the $W$ boson, we obtain a factor $i\frac{g_2}{\sqrt{2}}\gamma^\alpha P_L$. With this said, and using the Feynman rules we already know, we get
\[
\pazocal{M}(a,b,c,d)=\bar{u}_b(P_e)i\frac{g_2}{\sqrt{2}}\gamma^\beta P_Lv_c(P_1)\left[-i\frac{(g^{\alpha\beta}-k^\alpha k^\beta/m_W^2)}{K^2-m^2_W+i\varepsilon}\right]\bar{u}_d(P_2)i\frac{g_2}{\sqrt{2}}\gamma_\alpha P_Lu_a(P_\mu)
\]
At this point, it is time to use some approximations:
\[
\left\{
\begin{aligned}
&K^2\simeq P^2_\mu-2P_\mu\cdot P_2\sim\pazocal{O}(m_\mu^2)\\
&m_\mu\sim0.1\text{ GeV}\\
&m_W\sim10^2\text{ GeV}
\end{aligned}
\right.
\Rightarrow
\frac{i}{m^2_W}\left(\frac{ig_2}{\sqrt{2}}\right)^2=-\frac{i}{2}\underbrace{\left(\frac{g_2^2}{m_W^2}{\color{red}\frac{\sqrt{2}}{8}}\right)}_{G_F}{\color{red}\frac{8}{\sqrt{2}}}=-\frac{4iG_F}{\sqrt{2}}
\]
With these approximations, we get
\[
\pazocal{M}(a,b,c,d)=\bar{u}_b(P_e)\gamma^\beta\left(\frac{1-\gamma^5}{2}\right)v_c(P_1)\bar{u}_d(P_2)\gamma^\beta\left(\frac{1-\gamma^5}{2}\right)u_a(P_\mu)\left(-\frac{4iG_F}{\sqrt{2}}\right)
\]
What we are really interested in is $\frac{1}{2}\sum_{a,b,c,d}\pazocal{M}(a,b,c,d)\pazocal{M}^*(a,b,c,d)$, which can be expressed as
\[
\frac{1}{2}\sum_{a,b,c,d}\pazocal{M}(a,b,c,d)\pazocal{M}^*(a,b,c,d)=\left(\frac{G_F}{\sqrt{2}}\right)^2L_e^{\beta\gamma}(L_\mu)_{\beta\gamma}
\]
We now focus on $L_e^{\beta\gamma}$:
\begin{align*}
L_e^{\beta\gamma}&=\bar{u}_b(P_e)\gamma^\beta(1-\gamma^5)v_c(P_1)v_c^\dagger(P_1)(1-\gamma^5)^\dagger(\gamma^\gamma)^\dagger\bar{u}_b^\dagger(P_e)\\
&=\bar{u}_b(P_e)\gamma^\beta(1-\gamma^5)v_c(P_1)\bar{v}_c(P_1)\gamma^0(1-\gamma^5)\gamma^0\gamma^\gamma\gamma^0\gamma^0u_b(P_e)\\
&=\bar{u}_b(P_e)\gamma^\beta(1-\gamma^5)v_c(P_1)\bar{v}_c(P_1)(1+\gamma^5)\gamma^\gamma u_b(P_e)
\end{align*}
Now, keep in mind that $\sum u(p)\bar{u}(p)=\slashed{P}+m$ and $\sum v(p)\bar{v}(p)=\slashed{P}-m$. For our purpose, in this part we can set $m_e=m_{\bar{\nu}_e}=m_{\nu_\mu}=0$ and finally we get:
\[
L_e^{\beta\gamma}=4\Tr{\gamma^\beta\left(\frac{1-\gamma^5}{2}\right)\slashed{P_1}\left(\frac{1+\gamma^5}{2}\right)\gamma^\gamma\slashed{P_e}}=4\Tr{\gamma^\beta\left(\frac{1-\gamma^5}{2}\right)\slashed{P_1}\gamma^\gamma\slashed{P_e}}\marginnote{We used the fact that $P_L^2=P_L$}
\]
To calculate the trace of this object, we use the following rules:
\[
\left\{
\begin{aligned}
&\Tr{\gamma^\mu\gamma^\nu\gamma^\rho\gamma^\sigma}=4\left(g^{\mu\nu}g^{\rho\sigma}-g^{\mu\rho}g^{\nu\sigma}+g^{\mu\sigma}g^{\nu\rho}\right)\\
&\Tr{\gamma^\mu\gamma^\nu\gamma^\rho\gamma^\sigma\gamma^5}=-4i\varepsilon^{\mu\nu\rho\sigma}; \quad \varepsilon^{\alpha\beta\mu\nu}\varepsilon_{\alpha\beta\rho\sigma}=-2(\delta^\mu_\rho\delta^\nu_\sigma-\delta^\mu_\sigma\delta^\nu_\rho)
\end{aligned}
\right.
\]
After repeating everything for $(L_\mu)_{\beta\gamma}$ and tedious calculations, we obtain
\[
\frac{1}{2}\sum|\pazocal{M}|^2=64G^2_F(P_e\cdot P_2)(P_\mu\cdot P_1)
\]
We can compute the decay width, which is equal to
\[
\Gamma=\frac{1}{2m_\mu}\int\frac{d^3P_e}{(2\pi)^3}\frac{1}{2E_e}\int\frac{d^3P_1}{(2\pi)^3}\frac{1}{2E_1}\int\frac{d^3P_2}{(2\pi)^3}\frac{1}{2E_2}(2\pi)^4\delta^4(P_\mu-P_e-P_1-P_2)\langle|\pazocal{M}|^2\rangle
\]
Now, we integrate on $d^3P_e$:
\[
\int\frac{d^3P_e}{2E_e}=\int d^4P_e\theta(E_e)\delta(P_e^2-m_e^2)\delta^4(P_\mu-P_e-P_1-P_2)=\theta(E_\mu-E_1-E_2)\delta((P_\mu-P_1-P_2)^2)
\]
Suppose now that we fix the direction of $P_1$ as the x-axis of our system: we have no dependence on the azimuth angle because the 3-body decay is planar, so $\int d\Omega_1=4\pi$ and $\int d\phi_{12}=2\pi$. We are in a reference system such that $P_\mu=(m_\mu,\Vec{0})$.
\[
\Gamma=\frac{4\pi}{2m_\mu}\frac{2\pi}{(2\pi)^5}\int_0^{+\infty}\frac{dE_1E_1^2}{2E_1}\int_0^{+\infty}\frac{dE_2E_2^2}{2E_2}\int_{-1}^{+1}d\cos\theta_{12}\delta(m_\mu^2-2m_\mu(E_1+E_2)+2E_1E_2(1-\cos\theta_{12}))64G_F^2(P_e\cdot P_2)(P_\mu\cdot P_1)
\]
We integrate the delta function in $d\cos\theta_{12}$. We can rewrite it as:
\[
\delta\left(\cos\theta_{12}-\frac{m_\mu^2-2m_\mu(E_1+E_2)+2E_1E_2}{2E_1E_2}\right)
\]
Hence, we have to see for which values of $E_1$ and $E_2$, $\cos\theta_{12}$ is in the interval [-1,+1]:
\[
\left\{
\begin{aligned}
E_1&: 0\le E_1\le m_\mu/2\\
E_2&: -1\le\frac{m_\mu^2-2m_\mu(E_1+E_2)+2E_1E_2}{2E_1E_2}\le1\Rightarrow\frac{m_\mu}{2}-E_1\le E_2\le\frac{m_\mu}{2}
\end{aligned}
\right.
\]
The upper limit on $E_1$ corresponds to having $e$ and $\nu_\mu (2)$ emitted in one direction and $\bar{\nu}_e (1)$ emitted in the other one, hence splitting the muon energy $m_\mu$ equally in the to directions. For $E_2$, we fix $E_1$ and do tedious calculations. The final result is:
\[
\Gamma=\frac{1}{2^5m_\mu\pi^3}\int_0^{m_\mu/2}dE_1E_1\int_{m_\mu/2-E_1}^{m_\mu/2}dE_2E_22^6G_F^2(P_e\cdot P_2)(P_\mu\cdot P_1)=\frac{G_F^2m_\mu^5}{192\pi^3}=\Gamma_0
\]
If we include the effects of electrons, we get some corrections:
\[
\Gamma=\Gamma_0(1-8r+8r^3-r^4-12r^4\ln r) \quad r=(m_e/m_\mu)^2
\]
We can estimate this decay width by taking into account that\\
$G_F\sim1/300^4$ GeV$^{-4}$ and $m_\mu^5\sim10^{-5}$ GeV$^5$. It follows
\[
\Gamma\sim10^{-19}\text{ GeV}\Rightarrow\tau_\mu=2.1\cdot10^{-6}\text{ s}
\]


\section{How to build a model/Lagrangian/EFT?}
Our goal is to find the effective Lagrangian of nature, effective field theories (EFT) are characterised by certain degrees of freedom at a certain energy scale. To build a Lagrangian we have to:
\begin{enumerate}
    \item specify the local symmetries, i.e. the gauge group G of the theory+Lorentz
    \item specify the fields content
    \item write the most general Lagrangian according to 1. and 2.
    \item spend several hours admiring the beautiful Lagrangian we just wrote and try to understand it without summoning any ancient demon
\end{enumerate}
\textit{"Everything not forbidden is compulsory. -Marco Nardecchia"} -Michael Scott
\begin{figure}[h!]
\includegraphics{Images/WikiHow.png}
\caption{WikiHow is not a good place to learn how to build a Lagrangian.}
\end{figure}
\newline
Let's have a look at the case of QED:
\begin{table}[h]
    \centering
    \begin{tabular}{c|c|c}
        Field & G=U$(1)_C$ & Lorentz\\
        \hline
        $A_\mu$ & 0 & vector\\
        $\Psi_L$ & -1 & (1/2,0)\\
        $\Psi_R$ & -1 & (0,1/2)\\
        \hline
    \end{tabular}
    \caption*{}
    \label{tab:qed}
\end{table}
\newline
The most general Lagrangian up to dimension 4 is given by:
\[
\pazocal{L}=-\frac{1}{4}F^{\mu\nu}F_{\mu\nu}+i\bar{\Psi}\gamma^\mu D_\mu\Psi-m\bar{\Psi}\Psi+\text{h.c.}
\]
with $F_{\mu\nu}=\partial_\mu A_nu-\partial_\nu A_\mu$ and $D_\mu=i\partial_\mu-ieA_\mu$. Focus now on the Lagrangian without the mass term:
\[
\pazocal{L}=-\frac{1}{4}F^{\mu\nu}F_{\mu\nu}+i\bar{\Psi}_L\gamma^\mu D_\mu\Psi_L+i\bar{\Psi}_R\gamma^\mu D_\mu\Psi_R{\color{red}
-(m\bar{\Psi}_L\Psi_R+m^*\bar{\Psi}_R\Psi_L)}
\]
We can set $m=\hat{m}e^{i\theta}$ and $\Psi_L\xrightarrow[]{}\Psi'_L=\Psi_Le^{i\theta}$. In this way:
\[
(m\bar{\Psi}_L\Psi_R+m^*\bar{\Psi}_R\Psi_L)\xrightarrow[]{}(\hat{m}\bar{\Psi}'_L\Psi_R+\hat{m}\bar{\Psi}_R\Psi'_L)=m\bar{\Psi}\Psi
\]
In the case of QED, the group G=U$(1)_C=$U$(1)_L\times$U$(1)_R$ (2 generators) is broken into H=U$(1)_V$ (1 generator), i.e. $\Psi_{L/R}\xrightarrow[]{}e^{i\alpha}\Psi_{L/R}$:
\[
\left\{
\begin{aligned}
&\text{\# naive generators}=1 \text{ complex}=2 \text{ real}=\text{ modulus+phase}\\
&\text{\# broken generators}=2-1=1\\
&\text{\# physical parameters=\# naive-\# broken=1}
\end{aligned}
\right.
\]
\subsection{Review of the Standard Model}
For what concerns the Standard Model, we have:
\begin{enumerate}
    \item the gauge group is G=SU$(3)_C\times$SU$(2)_L\times$U$(1)_Y$
    \item the fields content is
    \begin{table}[h]
        \centering
        \begin{tabular}{c|c|c|c|c}
        Field & SU$(3)_C$ & SU$(2)_L$ & U$(1)_Y$ & Lorentz \\
        \hline
        $Q_L^i=\begin{pmatrix}u_L\\d_L
        \end{pmatrix}$ & 3 & 2 & 1/6 & (1/2,0) \\
        $u_R$ & 3 & 1 & 2/3 & (0,1/2) \\
        $d_R$ & 3 & 1 & -1/3 & (0,1/2) \\
        \hline
        $l_L^i=\begin{pmatrix}\nu_L\\e_L\end{pmatrix}$ & 1 & 2 & -1/2 & (1/2,0) \\
        $e_R$ & 1 & 1 & -1 & (0,1/2) \\
        \hline
        $H$ & 1 & 2 & 1/2 & \\
        $\Tilde{H}$ & 1 & 2 & -1/2 & \\
        \hline
        \end{tabular}
        \caption*{}
        \label{tab:my_label}
    \end{table}
    \item $\pazocal{L}=\pazocal{L}_{kin}+\pazocal{L}_\theta+\pazocal{L}_{Yuk}-V(H)$ with
    \[
    \left\{
    \begin{aligned}
    &\pazocal{L}_{kin}=\sum_{\{\Psi\}}\bar{\Psi}i\slashed{D}\Psi-\frac{1}{4}G_{\mu\nu}^aG^{\mu\nu,a}-\frac{1}{4}W_{\mu\nu}^iW^{\mu\nu,i}-\frac{1}{4}B_{\mu\nu}B^{\mu\nu}G_{\mu\nu}^aG^{\mu\nu,a}+(D_\mu H)^\dagger(D_\mu H) \\
    &D_\mu=\partial_\mu-ig_1YB_\mu-ig_2T^iW_\mu^i-ig_3\lambda^aG_\mu^a\\
    &\pazocal{L}_\theta=\frac{g_3^2}{32\pi^2}\theta G_{\mu\nu}^aG_{\rho\sigma}^a\varepsilon^{\mu\nu\rho\sigma} \\
    &\pazocal{L}_{Yuk}=-\sum_{i,j=1}^3\left[\bar{Q}_L^i(Y_u)^{ij}H^cu_R^i+\bar{Q}_L^i(Y_d)^{ij}Hd_r^j+\bar{l}_L^i(Y_e)^{ij}He_r^j+\text{h.c}\right] \quad H^c=i\sigma^2H^* \\
    &V(H)=-\mu^2H^\dagger H+\lambda(H^\dagger H)^2
    \end{aligned}
    \right.
    \]
\end{enumerate}
The mass terms are impossible to write because they are not gauge invariant:
\[
\left\{
\begin{aligned}
&\bar{Q}_LQ_L \quad &&\text{ not Lorentz invariant}\\
&\bar{Q}_Lu_R \quad &&\text{ not gauge invariant}\\
&\bar{Q}_Ld_R \quad &&\text{ not gauge invariant}
\end{aligned}
\right.
\]
The Higgs potential can be rewritten taking into account that $H$ is in the form
\[
H=\begin{pmatrix}\phi_1+i\phi_2\\\phi_3+i\phi_4\end{pmatrix}
\]
hence the term $H^\dagger H$ which appears in $V(H)$ is equal to $\varphi^2/2$:\marginnote{\includegraphics[]{Images/cowboy.jpg}\\
Do not wear a sombrero potential hat if you encounter a cowboy, that might be dangerous: cowboys can be evil.}
\[
V(H)\xrightarrow[]{}V(\phi)=-\frac{\mu^2}{2}\phi^2+\frac{\lambda}{4}\phi^4
\]
\begin{figure}[h!]
    \centering
    \includegraphics{Images/sombrero.png}
    \caption{Sombrero potential}
    \labfig{sombrero}
\end{figure}\\
This potential has two minima in $\pm v$, where $v=\sqrt{\mu^2/\lambda}=\langle\phi\rangle$: the field $\phi(x)$ can be written as its mean value plus an excitation which turns out to be the Higgs boson:\marginnote{\includegraphics[]{Images/higgs.jpg}\\
From \href{https://commons.wikimedia.org/wiki/File:Nobel_Prize_24_2013.jpg}{WikiMedia}: Peter Ware Higgs (born 29 May 1929) is a British theoretical physicist, Emeritus Professor in the University of Edinburgh, and Nobel Prize laureate for his work on the mass of subatomic particles.}
\[
\phi(x)=\frac{v+h(x)}{\sqrt{2}}
\]
We have 4 degrees of freedom in $H(x)$, so we can parameterize it as follows
\[
H(x)=e^{i\overset{\mathclap{\tikz \node {$\downarrow$} node [above=1ex] {\footnotesize  Nambu-Goldstone bosons};}}{\chi^i(x)}\sigma^i}\begin{pmatrix}0\\\frac{v+h(x)}{\sqrt{2}}\end{pmatrix}
\]
$h(x)$ is the radial excitation while $\chi(x)$ is the angular excitation: the latter costs no energy since $V(\phi)$ does not depend on $\chi$, therefore the Nambu-Goldstone bosons (NGB) are massless. The \href{https://en.wikipedia.org/wiki/Peter_Higgs}{Higgs} scalar potential breaks SU$(2)_{EW}\times$U$(1)_Y$ into U$(1)_Q$, i.e. 3+1 generators are broken into 1, resulting in 3 massless modes.

Now we bring our attention to the kinetic term, in particular to $|D_\mu H|^2$:
\begin{align*}
|D_\mu H|^2&=(D_\mu H)^\dagger(D_\mu H)=(\partial_\mu H)^\dagger(\partial_\mu H)+H^\dagger\left(\frac{g_1B_\mu}{\sqrt{2}}+g_2W_\mu^i\frac{\sigma^i}{2}\right)^2H\\
&+\frac{1}{2}H^\dagger(g_1B_\mu+\sigma^iW_\mu^ig_2)\partial^\mu H+\text{h.c.}\\
&=\frac{(v+h(x))^2}{8}\begin{pmatrix}0&1\end{pmatrix}\left[g_1B_\mu+g_2W_\mu^3\sigma^3+g_2(W_\mu^1\sigma^1+W_\mu^2\sigma^2)\right]^2\begin{pmatrix}0\\1\end{pmatrix}\\
&=(v+h(x))^2\left[\frac{(g_1^2+g_2^2)}{8}Z_\mu Z^\mu+\frac{g_2^2}{4}W_\mu^+W^{\mu-}\right]
\end{align*}
where we introduced $W_\mu^\pm=(W_\mu^1\pm iW_\mu^2)/\sqrt{2}$. Moreover, we can define $m_Z$ and $m_W$ as:
\[
\left\{
\begin{aligned}
m_Z&=\frac{\sqrt{g_1^2+g_2^2}}{2}v\\
m_W&=\frac{g_2}{2}v
\end{aligned}
\right.
\Rightarrow|D_\mu H|^2=\left(1+\frac{h(x)}{v}\right)^2\left[\frac{m_Z^2}{2}Z_\mu Z^\mu+m_W^2W_\mu^+W^{\mu-}\right]
\]
Notice that
\[
\frac{m_W^2}{m_Z^2}=\frac{g_2^2}{g_1^2+g_2^2}=\cos^2\theta_W
\]
from which we can define the quantity $\rho$: 
\[
\rho=\frac{m_W^2}{m_Z^2\cos^2\theta_W}
\]
which is equal to 1 at tree level but this symmetry gets broken by the hypercharge.

The field $H(x)$ can be written in such a way that there is mixing between $\phi_i$ and the gauge fields:
\[
H(x)=\begin{pmatrix}\phi_1+i\phi_2\\\frac{v+h(x)+i\phi_3}{\sqrt{2}}\end{pmatrix}
\]
We define two new matrices $\sigma^+$ and $\sigma^-$ and redo all the computations for $|D_\mu H|^2$
\[
\sigma^+=\begin{pmatrix}0&1\\0&0\end{pmatrix}
\quad
\sigma^-=\begin{pmatrix}0&0\\1&0\end{pmatrix}
\]
\begin{align*}
|D_\mu H|^2&=iH^\dagger[g_1B_\mu+g_2\sigma_3W_\mu^3+g_2(W_\mu^1\sigma^1+W_\mu^2\sigma^2)]\partial_\mu H+\text{h.c.}\\
&=\frac{v+h(x)}{\sqrt{2}}\begin{pmatrix}0&1\end{pmatrix}\left[(g_1B_\mu-g_2W_\mu^3)\frac{1}{\sqrt{2}}\partial_\mu(h+i\phi_3)\begin{pmatrix}0\\1\end{pmatrix}+\text{h.c.}\right]\\
&+\frac{v+h(x)}{\sqrt{2}}\begin{pmatrix}0&1\end{pmatrix}\left[\frac{g_2}{\sqrt{2}}(W_\mu^+\sigma^++W_\mu^-\sigma^-)\partial_\mu H+\text{h.c.}+\dots\right]\\
&=\frac{v+h(x)}{2}\left[-(g_1B_\mu-g_2W_\mu^3)\partial_\mu\phi_3+(ig_2W_\mu^-\partial_\mu\phi^++\text{h.c.})+\dots\right]\\
&=m_Z\partial_\mu\phi_3Z^\mu+(im_WW_\mu^-\partial_\mu\phi^++\text{h.c.})+\dots
\end{align*}
We now add a gauge-fixing term to the Lagrangian:
\[
\pazocal{L}_{GF}=-\frac{1}{2\xi}(\partial_\mu A^\mu)^2-\frac{1}{2\xi_Z}(\partial_\mu Z^\mu+m_Z\phi_3\xi_Z)^2-\frac{1}{2\xi_W}|\partial_\mu W^{-\mu}-im_W\xi_W\phi^-|^2
\]
For the $Z$ boson, the propagator is given by:
\[
-\frac{i}{q^2-m_z^2}\left[\eta_{\mu\nu}-(1-\xi_Z)\frac{q_\mu q_\nu}{q^2-\xi_Zm_Z^2}\right]
\]
In the unitary gauge, we have that $\xi_{W/Z}\to\infty$ hence $\phi_3$ has to vanish while in the Landau gauge $\xi_{W/Z}\to0$: $\partial_\mu Z^\mu$ has to vanish.

To express $D_\mu$ in terms of the physical fields, we have to perform a rotation:
\[
\left(\begin{array}{c}
      B_\mu\\
      W_\mu^3
\end{array}\right)=\left(\begin{array}{cc}
    \cos\theta_W & \sin\theta_W \\
    -\sin\theta_W & \cos\theta_W
\end{array}\right)\left(\begin{array}{c}
      A_\mu\\
      Z_\mu
\end{array}\right)
\]
After this rotation, we can write:
\[
D_\mu=\partial_\mu-i\underset{\mathclap{\tikz \node {$\uparrow$} node [below=1ex] {\footnotesize  em coupling: $1/e^2=1/g_1^2+1/g_2^2$};}}{e}\overset{\mathclap{\tikz \node {$\downarrow$} node [above=1ex] {\footnotesize  generator $Q=T_3+Y$};}}{Q}A_\mu-i\frac{e}{\cos\theta_W\sin\theta_W}(T_3-\sin^2\theta_WQ)Z_\mu-i\frac{g_2}{\sqrt{2}}(\underset{\mathclap{\tikz \node {$\uparrow$} node [below=1ex] {\footnotesize  $T^\pm=\sigma^\pm/2$};}}{T^+}W_\mu^++T^-W_\mu^-)-ig_3\overset{\mathclap{\tikz \node {$\downarrow$} node [above=1ex] {\footnotesize  Gell-Mann matrices};}}{\lambda^a}G_\mu^a
\]
At this point, we move our attention to the Yukawa sector: we know that it is written as
\[
\pazocal{L}_{Yuk}=-\bar{Q}_L^i(Y_u)^{ij}H^cu_R^i+\bar{Q}_L^i(Y_d)^{ij}Hd_r^j+\bar{l}_L^i(Y_e)^{ij}He_r^j+\text{h.c}
\]
In the unitary gauge, we have $H=\begin{pmatrix}0\\v/\sqrt{2}\end{pmatrix}\begin{pmatrix}1+h/v\end{pmatrix}$ and the Yukawa interaction can be rewritten in the following form:
\[
\pazocal{L}_{Yuk}\supset-\bar{u}_L^i\left(Y_u^{ij}\frac{v}{\sqrt{2}}\right)u_R^j\left(1+\frac{h}{v}\right)-\bar{d}_L^i\left(Y_d^{ij}\frac{v}{\sqrt{2}}\right)d_R^j\left(1+\frac{h}{v}\right)-\bar{e}_L^i\left(Y_e^{ij}\frac{v}{\sqrt{2}}\right)l_R^j\left(1+\frac{h}{v}\right)+\text{h.c}
\]
Furthermore, we can rotate up left/right and down left/right:\marginnote{\includegraphics{Images/cabibbo.jpeg}\\
From \href{https://www.nature.com/articles/nphys1791}{Nature}: Italian particle physicist Nicola Cabibbo (10 April 1935-16 August 2010) had been a towering figure in Italian physics, particularly as president of the National Institute of Nuclear Physics (INFN) between 1983 and 1992, and then as president of the Pontifical Academy of Sciences.}
\[
\left\{
\begin{aligned}
&u_L^i\xrightarrow[]{}U_L^{ij}u_L^j\\
&u_R^i\xrightarrow[]{}U_R^{ij}u_R^j
\end{aligned}
\right.
\quad
\left\{
\begin{aligned}
&d_L^i\xrightarrow[]{}D_L^{ij}d_L^j\\
&d_R^i\xrightarrow[]{}D_R^{ij}d_R^j
\end{aligned}
\right.
\quad
U_{L/R},D_{L/R}\in\text{SU}(3)
\]
We introduce the \href{https://en.wikipedia.org/wiki/Cabibbo-Kobayashi-Maskawa_matrix}{Cabibbo–Kobayashi–Maskawa (CKM) matrix}\\ $V_{CKM}^{ij}=(U_LD_L^\dagger)^{ij}$ which is a special unitary $3\times3$ matrix with 8 real parameters: 3 angles+5 phases, with 4 of them being unphysical. The interaction with the Higgs implies flavour mixing but there is no flavour change at tree level.
\section{Symmetries}
\subsection{Parity}
We already know that we can parameterize our field as composed by a left-handed spinor $\xi$ and a right-handed spinor $\eta$:
\[
\Psi=\left(\begin{array}{c}
     \xi\\
     \eta
\end{array}\right)
\quad
\left\{
\begin{aligned}
&\xi=(2,1) \text{ of Lorentz}\\
&\eta=(1,2) \text{ of Lorentz}
\end{aligned}
\right.
\]
Moreover, we have that under parity position and momentum transform as:
\[
(\Vec{x},t)\xrightarrow[]{}(-\Vec{x},t) \quad \Vec{p}\xrightarrow[]{}-\Vec{p}
\]
while the spin remains untouched being a pseudo-vector. For the right- and left-handed spinors, we have the relation $\xi(\Vec{x,t})\xleftrightarrow[]{}\eta(-\Vec{x},t)$ and we can introduce an intrinsic parity of the field, i.e. a phase:
\[
\left\{
\begin{aligned}
&\xi(\Vec{x},t)\xrightarrow[]{}\zeta_p\eta(-\Vec{x},t)\\
&\eta(\Vec{x},t)\xrightarrow[]{}\zeta_p\xi(-\Vec{x},t)
\end{aligned}
\right.
\]
In terms of $\Psi$ we have $\Psi(\Vec{x},t)\xrightarrow[]{}\zeta_p\gamma^0\Psi(-\Vec{x},t).$\marginnote{The gamma matrices are defined as:
\[
\gamma^\mu=\left(\begin{array}{cc}
    0 & \sigma^\mu \\
    \Bar{\sigma}^\mu & 0
\end{array}\right), \gamma^5=\left(\begin{array}{cc}
    -1 & \\
     & -1
\end{array}\right)
\]
with $\sigma^\mu=(\mathbb{1},\sigma^i)$, $\bar{\sigma}^\mu=(\mathbb{1},-\sigma^i)=\eta_{\mu\mu}\sigma^\mu$ and $\eta_{\mu\mu}=(+1,-1,-1,-1)$} We can have a look at the current under parity:
\begin{align*}
\bar{\Psi}\gamma^\mu\Psi=\Psi^\dagger\gamma^0\gamma^\mu\Psi=\xi^\dagger\bar{\sigma}^\mu\xi+\eta^\dagger\sigma^\mu\eta\xrightarrow[\text{parity}]{}&\eta^\dagger\bar{\sigma}^\mu\eta+\xi^\dagger\sigma^\mu\xi\\
&=\eta_{\mu\mu}(\eta^\dagger\sigma^\mu\eta+\xi^\dagger\bar{\sigma}^\mu\xi)\\
&=\eta_{\mu\mu}\bar{\Psi}\gamma^\mu\Psi
\end{align*}
This is just a vector, now we include $A_\mu(\Vec{x},t)\xrightarrow[]{}\eta_{\mu\mu}A_\mu(-\Vec{x},t)$: it follows that $A_\mu\bar{\Psi}\gamma^\mu\Psi$ is invariant, hence QED is invariant under parity. The same results apply for QCD since $G_\mu^a\bar{\Psi}\lambda^a\gamma^\mu\Psi$ is invariant. On the other hand, the theta term is odd under parity: it is proportional to $G_{\mu\nu}G_{\rho\sigma}\varepsilon^{\mu\nu\rho\sigma}$ and $\varepsilon$ forces 3 indices to be space-like and 1 to be time-like.
\subsection{Charge conjugation}
We now see what happens when we transform the field $\Psi$ under charge conjugation:
\[
\Psi(x)\xrightarrow[]{C}i\gamma_2\Psi^*(x)\Tilde{\eta}_\Psi
\quad
\left\{
\begin{aligned}
&\xi(x)\xrightarrow[]{C}\Tilde{\eta}_\Psi\eta^c(\Vec{x},t) \,\text{ with }\, \eta^c=i\sigma^2\eta^*\\
&\eta(x)\xrightarrow[]{C}-\Tilde{\eta}_\Psi\xi^c(\Vec{x},t) \,\text{ with }\, \xi^c=i\sigma^2\xi^*
\end{aligned}
\right.
\]
Moreover, we have that $\eta^\dagger\sigma^\mu\eta=-\eta^c^\dagger\bar{\sigma}^\mu\eta^c$:
\[
-\eta^c^\dagger\bar{\sigma}^\mu\eta^c=-(i\sigma^2\eta^*)^\dagger\bar{\sigma}^\mu(i\sigma^2\eta^*)=-\eta^T\underbrace{\sigma_2\bar{\sigma}^\mu\sigma_2}_{\sigma^\mu^T}\eta^*=-\eta^T\sigma^\mu^T\eta^*=\eta^\dagger\sigma^\mu\eta
\]
What happens now if we do $C^\dagger A_\mu(x)C$? We obtain just a phase $\eta_A$ since photons are real particles.

The QCD case is more complicated: we have $\bar{\Psi}_\alpha\gamma^\mu\frac{\lambda^a_{\alpha\beta}}{2}G_\mu^a\Psi_\beta$ and we ask ourselves how does $\bar{\Psi}_\alpha\gamma^\mu\Psi_\beta$ transform under charge conjugation.
\[
\left\{
\begin{aligned}
&\bar{\Psi}_\alpha\gamma^\mu\Psi_\beta\xrightarrow[]{C}-\bar{\Psi}_\beta\gamma^\mu\Psi_\alpha\Tilde{\eta}_{\Psi_\alpha}\Tilde{\eta}^*_{\Psi_\beta}\\
&\bar{\Psi}_\alpha\gamma^\mu\frac{\lambda^a_{\alpha\beta}}{2}\Psi_\beta\xrightarrow[]{C}\bar{\Psi}\gamma^\mu\frac{(\lambda^a)^T}{2}\Psi\\
&G_\mu^a\lambda^a\xrightarrow[]{C}-G_\mu^a(\lambda^a)^T\Rightarrow C^\dagger G_\mu^a C=\Tilde{\eta}_{G_\mu}G_\mu^a
\end{aligned}
\right.
\]
The phase $\Tilde{\eta}_{G_\mu}$ depends on the generators:
\[
\left\{
\begin{aligned}
&(\lambda^a)^T=\lambda^a &&\text{ for } \, a=1,3,4,6,8\Rightarrow\Tilde{\eta}_{G_\mu}=-1\\
&(\lambda^a)^T=-\lambda^a &&\text{ for } \, a=2,5,7\Rightarrow\Tilde{\eta}_{G_\mu}=+1
\end{aligned}
\right.
\]
\subsection{Discrete Symmetries}
We are now going to review discrete symmetries, focusing on states, creation/annihilation  operators and fields:\marginnote{$A$ is just an internal quantum number.}
\[
\ket{A,\Vec{p},\sigma}\xrightarrow[\text{parity}]{}\left\{
\begin{aligned}
&\hat{P}\ket{A,\Vec{p},\sigma}=\eta_A\ket{A,-\Vec{p},\sigma}\\
&\hat{P}\ket{\bar{A},\Vec{p},\sigma}=\eta_{\bar{A}}\ket{\bar{A},-\Vec{p},\sigma}
\end{aligned}
\right.
\]
The operators $\hat{a}^\dagger, \hat{a}$ refer to particles while $\hat{b}^\dagger, \hat{b}$ refer to anti-particles:
\[
\hat{P}\hat{a}^\dagger_\sigma(\Vec{p})\hat{P}^\dagger\ket{0}=\eta_A\hat{a}^\dagger_\sigma(-\Vec{p})\ket{0}\Rightarrow\left\{
\begin{aligned}
&\hat{P}\hat{a}^\dagger_\sigma(\Vec{p})\hat{P}^\dagger=\eta_A\hat{a}^\dagger_\sigma(-\Vec{p})\\
&\hat{P}\hat{b}^\dagger_\sigma(\Vec{p})\hat{P}^\dagger=\eta_{\bar{A}}\hat{b}^\dagger_\sigma(-\Vec{p})
\end{aligned}
\right.
\]
Now we want the transformation rules for the fields: $\hat{P}\phi(x)\hat{P}^\dagger=?$ For scalar fields, we have:
\[
\hat{P}\phi(x)\hat{P}^\dagger=\eta_\phi\phi(\pazocal{P}x) \, \text{ with } \, \pazocal{P}=\text{diag}(+1,-1,-1,-1)
\]
Furthermore, we know how to express the field $\phi(x)$ in terms of creation and annihilation operators:
\[
\phi(x)=\int \frac{d^3p}{(2\pi)^3}\frac{1}{\sqrt{2E_p}}\left[\hat{a}(\Vec{p})e^{-ipx}+\hat{b}^\dagger(\Vec{p})e^{+ipx}\right]
\]
Acting with parity on $\phi(x)$, in order to get the previous form, it has to be $\eta_\phi=\eta_A^*$ and $\eta_A\eta_{\bar{A}}=1$.

It is possible to repeat the same procedure for fermions:
\[
\Psi(x)=\sum_\sigma\int \frac{d^3p}{(2\pi)^3}\frac{1}{2E(\Vec{p})}\left[\hat{a}_\sigma(\Vec{p})e^{-ipx}u(\Vec{p})+\hat{b}^\dagger_\sigma(\Vec{p})e^{+ipx}v(\Vec{p})\right]
\]
We now act with the parity operator $\hat{P}$:
\[
\hat{P}\Psi(x)\hat{P}^\dagger=\eta_\Psi\gamma^0\Psi(\pazocal{P}x)=\sum_\sigma\int \frac{d^3p}{(2\pi)^3}\frac{1}{2E(\Vec{p})}\left[\eta_A^*\hat{a}_\sigma(-\Vec{p})e^{-ipx}u(\Vec{p})+\eta_{\bar{A}}\hat{b}^\dagger_\sigma(-\Vec{p})e^{+ipx}v(\Vec{p})\right]
\]
The field $\Psi$ is now not calculated in $x$ but in $\pazocal{P}x$. For the spinors, this means $u(\Vec{p})=\gamma^0u(\pazocal{P}\Vec{p})$ and $v(\Vec{p})=-\gamma^0v(\pazocal{P}\Vec{p})$. What we get at the end is that $\eta_\Psi=\eta_A^*$ and $\eta_A\eta_{\bar{A}}=-1$.

Suppose now to have an Abelian symmetry $U(\theta)\ket{A,\Vec{p},\sigma}=e^{iq_A\theta}\ket{A,\Vec{p},\sigma}$ and to apply parity to it:
\[
\hat{P}U(\theta)\ket{A,\Vec{p},\sigma}=\underbrace{e^{iq_A\theta}\eta_A}_{\eta'_A}\ket{A,-\Vec{p},\sigma}
\]
We can redefine $\eta'_A$ to be equal to 1. Doing this, we also define a new operator $\hat{P}'=\hat{P}U(\theta=-\arg(\eta_A)/q_A)$. This is possible only when there is a global symmetry associated.
\begin{table}[h]
    \centering
    \begin{tabular}{c|c|c|c}
    State & $\eta$ & Global symmetry & Theory\\
    \hline
    Lepton & +1 & U$(1)_e\times$U$(1)_\mu\times$U$(1)_\tau$ & QED\\
    Anti-lepton & -1 & U$(1)_e\times$U$(1)_\mu\times$U$(1)_\tau$ & QED\\
    Quarks & +1 & U$(1)^{N_F}$ & QCD+QED\\
    Anti-quarks & -1 & U$(1)_u\times$U$(1)_d\times$U$(1)_s$ & QCD+QED\\
    \hline
    \end{tabular}
    \caption*{}
    \label{tab:my_label}
\end{table}
\end{document}